\documentclass[12pt]{article}
\usepackage[utf8]{inputenc}
\usepackage[T1]{fontenc}
\usepackage{hyperref}
\title{Wasa Language Resources for K-12 Curriculum Development}
\author{Compiled by Claude Sonnet 5.5 (claude-sonnet-5-5) for Andr\'as Kornai}
\date{}
\begin{document}
\maketitle

\section{Introduction}
Wasa (also Wassa or Wasaw; ISO 639-3 wss; Glottocode wasa1244) is an Akan language of southwestern Ghana, closely related to Twi and Fante, with a few speakers in C\^ote d'Ivoire. Speaker estimates are about 270,000 (Wikipedia, 2013) and 309,000 (German Wikipedia, 2003). Its dialects are Amenfi and Fianse. The Wasa people number about 382,000 according to Joshua Project, a separate figure from the language count. Because Wasa is close to the other Akan varieties, schooling in Ghana is likely to use those standardized varieties rather than Wasa itself.

\section{Primary Educational Resources}
No Wasa-specific primers, textbooks or school programmes were found.

\section{Linguistic and Reference Materials}
\begin{itemize}
\item Wikipedia on the Wasa language: \url{https://en.wikipedia.org/wiki/Wasa_language}
\item Wikipedia on the Wasa people: \url{https://en.wikipedia.org/wiki/Wasa_people}
\item German Wikipedia entry: \url{https://de.Wikipedia.Org/wiki/Wasa_(Sprache)}
\end{itemize}

\section{Religious and Community Resources}
\begin{itemize}
\item Joshua Project profile for the Wasa in Ghana: \url{https://joshuaproject.net/index.php/people_groups/15878}
\end{itemize}

\section{Digital and Audio Resources}
Global Recordings Network hosts audio Bible materials in Wassa: \url{https://globalrecordings.net/en/language/wss}

\section{Curriculum Development Context}
Mutual intelligibility with Fante, Akuapem, Asante and Abron means existing Akan educational materials could be adapted. That adaptation would need local community input on the Wasa-specific features. No grammar, dictionary or Wasa orthography was found in the search results.

\end{document}